%! AP Statistics — Unit 4, Lesson 4.6 — Exit Ticket KEY
\documentclass[10pt]{article}
\usepackage{apstats-article}
\usepackage{apstats-key}

\begin{document}

\pageheader{Unit 4, Lesson 4.6}{Exit Ticket --- KEY}
\namedateperiod

\vspace{0.14in}

\begin{scenariobox}[Context]{sky}
\small
School-week preference $\times$ grade level survey (160 students).

\vspace{6pt}
\renewcommand{\arraystretch}{1.4}
\begin{tabularx}{\linewidth}{@{} l X X r @{}}
 & \textbf{Underclassmen} & \textbf{Upperclassmen} & \textbf{Total} \\
\hline
\textbf{Prefers 4-day}  & 40 & 60 & 100 \\
\textbf{Prefers 5-day}  & 36 & 24 & 60  \\
\hline
\textbf{Total}          & 76 & 84 & 160 \\
\end{tabularx}
\end{scenariobox}

\vspace{0.1in}

\begin{practicebox}
\small
\begin{enumerate}[leftmargin=1.5em, itemsep=12pt]
  \item \ans{$P(\text{Prefers 4-day}) = 100/160 = 0.625$; $P(\text{Prefers 4-day}|\text{Upper}) = 60/84 \approx 0.714$.} \ans{Since $0.625 \neq 0.714$, Prefers 4-day week and Upperclassmen are \textbf{not independent} --- upperclassmen are more likely to prefer a 4-day week than the overall student population.}

  \item \ans{$100/160 + 84/160 - 60/160 = (100 + 84 - 60)/160 = 124/160 = 0.775$}

  \item \ans{The student has the two concepts confused. Mutually exclusive means $P(A \cap B) = 0$ --- a student cannot prefer both a 4-day and a 5-day week simultaneously. However, mutually exclusive events with positive probabilities are \emph{dependent}, not independent. If you know a student prefers the 4-day week, the probability they prefer the 5-day week drops to 0, which differs from $P(\text{5-day}) = 60/160 = 0.375$. Therefore the two events are dependent, not independent.}
\end{enumerate}
\end{practicebox}

\end{document}
